% *** ก่อนส่ง: ลบข้อความตัวอย่าง (\ldots) ทั้งหมด แล้วเขียนเนื้อหาของตนเอง ***
% =============================================================
% บทที่ 4 - การวิเคราะห์ระบบ และพัฒนาโปรแกรม
% แก้ไขเนื้อหาในบทนี้ได้ที่ไฟล์ chapter4.tex
% =============================================================

% -----------------------------------------------------------
\section{การวิเคราะห์ระบบ}
% -----------------------------------------------------------

\subsection{ผังแสดงทิศทางของกิจกรรม}

จากการวิเคราะห์ปัญหา สามารถแสดงผังทิศทางกิจกรรมของระบบได้ดังนี้\ldots{}

\subsection{รายการความต้องการ}

\subsubsection{ความต้องการเชิงฟังก์ชัน}

ความต้องการเชิงฟังก์ชันของระบบแสดงดังตารางที่ \ref{tab:func_req}

\begin{table}[H]
\centering
\caption{ความต้องการเชิงฟังก์ชัน (Functional Requirements)}
\label{tab:func_req}
\begin{tabular}{|c|l|l|}
\hline
\textbf{รหัส} & \textbf{ชื่อ Requirement} & \textbf{รายละเอียด} \\
\hline
FR-01 & \ldots & ระบบต้องสามารถ\ldots \\ \hline
FR-02 & \ldots & ระบบต้องสามารถ\ldots \\ \hline
FR-03 & \ldots & ระบบต้องสามารถ\ldots \\
\hline
\end{tabular}
\end{table}

\subsubsection{ความต้องการที่ไม่ใช่เชิงฟังก์ชัน}

ระบบต้องรองรับ\ldots{}

\begin{enumerate}[label=\arabic*)]
    \item \textbf{ด้านประสิทธิภาพ (Performance):} ระบบต้องตอบสนองภายใน\ldots{} วินาที
    \item \textbf{ด้านความปลอดภัย (Security):} ระบบต้องเข้ารหัสรหัสผ่านด้วย\ldots{}
    \item \textbf{ด้านความสามารถใช้งาน (Usability):} ระบบต้องใช้งานง่ายโดย\ldots{}
\end{enumerate}

\subsection{Use Case Diagram}

จากการวิเคราะห์ความต้องการ สามารถสร้าง Use Case Diagram ได้ดังภาพที่ \ref{fig:usecase}

% *** uncomment ด้านล่างเมื่อมีภาพ Use Case แล้ว ***
% *** เปลี่ยนชื่อไฟล์จาก usecase.png เป็นชื่อไฟล์ของคุณ ***
\begin{figure}[H]
\centering
% \includegraphics[width=0.9\textwidth]{images/usecase.png}
\fbox{\parbox{0.8\textwidth}{\centering\vspace{3cm}\textit{แทรกภาพ Use Case Diagram ที่นี่}\\ใช้คำสั่ง \texttt{\textbackslash includegraphics}\vspace{3cm}}}
\caption{Use Case Diagram ของระบบ}
\label{fig:usecase}
\end{figure}

\subsection{System Sequence Diagram}

ภาพรวมการทำงานของระบบแสดงด้วย System Sequence Diagram ดังภาพที่\ldots{}

\subsection{Conceptual Class Diagram}

ความสัมพันธ์เชิงแนวคิดของคลาสในระบบแสดงดังภาพที่\ldots{}

% -----------------------------------------------------------
\section{การออกแบบระบบ}
% -----------------------------------------------------------

\subsection{สถาปัตยกรรมของระบบ}

ระบบออกแบบด้วยสถาปัตยกรรมแบบ\ldots{} แสดงด้วย Package Diagram ดังภาพที่\ldots{}

\subsection{Class Diagram}

รายละเอียดคลาสและความสัมพันธ์แสดงดังภาพที่\ldots{}

\subsection{Sequence Diagram}

พฤติกรรมเชิงพลวัตของระบบแสดงด้วย Sequence Diagram ดังภาพที่\ldots{}

\subsection{การออกแบบฐานข้อมูล}

ออกแบบฐานข้อมูลเชิงสัมพันธ์ โดยแสดง E-R Diagram ดังภาพที่\ldots{}

% ตัวอย่างตารางโครงสร้างฐานข้อมูล
รายละเอียดโครงสร้างของตารางแต่ละตารางแสดงดังตารางที่ \ref{tab:db_users}

\begin{table}[H]
\centering
\caption{โครงสร้างตาราง users}
\label{tab:db_users}
\begin{tabular}{|l|l|l|l|}
\hline
\textbf{ฟิลด์} & \textbf{ชนิดข้อมูล} & \textbf{คีย์} & \textbf{คำอธิบาย} \\
\hline
id & INT & PK & รหัสผู้ใช้ \\ \hline
username & VARCHAR(50) & UQ & ชื่อผู้ใช้ \\ \hline
email & VARCHAR(100) & & อีเมล \\ \hline
password & VARCHAR(255) & & รหัสผ่าน (hashed) \\ \hline
created\_at & DATETIME & & วันที่สร้าง \\
\hline
\end{tabular}
\end{table}

\subsection{Mockup User Interface}

ต้นแบบส่วนติดต่อผู้ใช้ออกแบบดังภาพที่\ldots{}

% -----------------------------------------------------------
\section{การพัฒนาโปรแกรม}
% -----------------------------------------------------------

\subsection{ภาษาและเครื่องมือที่ใช้ในการพัฒนา}

ระบบพัฒนาด้วยภาษา\ldots{} ร่วมกับเฟรมเวิร์ค\ldots{} ใช้ฐานข้อมูล\ldots{}

\subsection{โครงสร้างของโปรแกรม}

โครงสร้างโปรแกรมแบ่งออกเป็น\ldots{} ส่วน ดังนี้\ldots{}

\subsection{ข้อกำหนดเฉพาะของแต่ละส่วน}

รายละเอียดข้อกำหนดของแต่ละส่วน (Program Specification) มีดังนี้\ldots{}

\subsection{การติดต่อกับฐานข้อมูลและ Services}

ระบบติดต่อกับฐานข้อมูลผ่าน\ldots{}

% -----------------------------------------------------------
\section{การทดสอบระบบ}
% -----------------------------------------------------------

\subsection{วิธีการทดสอบ}

ทดสอบด้วยวิธี User Acceptance Testing โดยมี Test Case ดังตารางที่ \ref{tab:test_result}

\subsection{ผลการทดสอบ}

ผลการทดสอบระบบแสดงดังตารางที่ \ref{tab:test_result}

\begin{table}[H]
\centering
\caption{ผลการทดสอบระบบ (User Acceptance Testing)}
\label{tab:test_result}
\begin{tabular}{|c|l|l|c|}
\hline
\textbf{TC} & \textbf{รายการทดสอบ} & \textbf{ผลที่คาดหวัง} & \textbf{ผล} \\
\hline
TC-01 & เข้าสู่ระบบด้วยข้อมูลถูกต้อง & เข้าสู่หน้าหลักได้ & ผ่าน \\ \hline
TC-02 & เข้าสู่ระบบด้วยรหัสผ่านผิด & แสดง error message & ผ่าน \\ \hline
TC-03 & \ldots & \ldots & ผ่าน \\ \hline
TC-04 & \ldots & \ldots & ผ่าน \\ \hline
TC-05 & \ldots & \ldots & ผ่าน \\
\hline
\end{tabular}
\end{table}

% ถ้าตารางยาวเกิน 1 หน้า ใช้ \captioncont ดังนี้:
% \begin{table}[H]
% \centering
% \captioncont{ผลการทดสอบระบบ (User Acceptance Testing)}
% \begin{tabular}{|c|l|l|c|}
% ...ข้อมูลส่วนที่เหลือ...
% \end{tabular}
% \end{table}

\subsection{รายละเอียดกรณีทดสอบ (Test Case Detail)}

% ตัวอย่าง Test Case แบบละเอียดพร้อมขั้นตอน
% *** คัดลอก template ด้านล่างไปใช้ แล้วเปลี่ยนข้อมูล ***
รายละเอียดกรณีทดสอบ TC-01 แสดงดังตารางที่ \ref{tab:tc01}

\begin{table}[H]
\caption{กรณีทดสอบ TC-01: เข้าสู่ระบบด้วยข้อมูลถูกต้อง}
\label{tab:tc01}
\noindent
\begin{tabular}{|l|l|}
\hline
\textbf{Test Case ID} & TC-01 \\ \hline
\textbf{ชื่อกรณีทดสอบ} & เข้าสู่ระบบด้วยข้อมูลถูกต้อง \\ \hline
\textbf{วัตถุประสงค์} & ทดสอบว่าผู้ใช้สามารถเข้าสู่ระบบได้เมื่อกรอกข้อมูลถูกต้อง \\ \hline
\textbf{เงื่อนไขก่อนทดสอบ} & มี account ในระบบอยู่แล้ว \\ \hline
\textbf{ขั้นตอนการทดสอบ} & 1. เปิดหน้า Login \\ \cline{2-2}
 & 2. กรอก Username: \texttt{testuser} \\ \cline{2-2}
 & 3. กรอก Password: \texttt{P@ssw0rd} \\ \cline{2-2}
 & 4. กดปุ่ม ``เข้าสู่ระบบ'' \\ \hline
\textbf{ผลที่คาดหวัง} & ระบบแสดงหน้าหลัก (Dashboard) พร้อมชื่อผู้ใช้ \\ \hline
\textbf{ผลการทดสอบจริง} & ระบบแสดงหน้า Dashboard ถูกต้อง \\ \hline
\textbf{สถานะ} & \textbf{ผ่าน} \\
\hline
\end{tabular}
\end{table}

